\documentclass[10pt,a4paper]{amsart}
\usepackage[margin=19mm]{geometry}
\usepackage{amsmath,amssymb,mathtools}
\usepackage[scr=rsfs,cal=euler]{mathalfa}
\usepackage{xfrac}
\usepackage[unicode,hidelinks,bookmarks=false]{hyperref}
\newcommand{\ie}{i.e.\ }
\newcommand{\cf}{cf.\ }
\newcommand{\resp}{respectively\ }
\newcommand{\Subsection}{\subsection}
\providecommand{\nobreakdash}{\nobreak\hbox{-}}
\setlength{\emergencystretch}{2em}
\multlinegap0pt
\usepackage{amssymb}
\usepackage{mathtools}
\usepackage[shortlabels]{enumitem}
\newtheorem{theorem}{Theorem}
\newcommand{\CC}{\mathcal{CC}}
\newcommand{\dotcupunion}{\mathbin{\dot\cup}}
\title{Clique decompositions and covers for large graphs}
\author{Wentao Zhang}
\date{Source: arXiv:2608.25233v1 (CC BY 4.0)}
\begin{document}
\maketitle

\noindent\emph{Source and licence.} Clique decompositions and covers for large graphs. Version arXiv:2608.25233v1 (CC BY 4.0), distributed by arXiv under the Creative Commons Attribution 4.0 International licence, http://creativecommons.org/licenses/by/4.0/, as recorded in the arXiv licence metadata for this exact version and shown on the abstract page https://arxiv.org/abs/2608.25233v1. Full licence notice: \texttt{licenses/arxiv-2608.25233-CC-BY-4.0.txt}.\par\smallskip

\noindent\textit{Editorial scope.} Counted unit: the article's theorem cov:thm:transversal-matching (source0.tex lines 2105--2126). The article's complete original proof of it is delivered with it (source0.tex lines 2128--2283, one \texttt{proof} environment of 156 lines). No mathematics is written, repaired or re-derived: the statement and the proof are the article's own lines, in the article's own notation, and no retained line is substituted or re-flowed.  No further source-local context is needed: every cross-reference of every retained line resolves inside the two delivered spans.  Nothing was omitted from any retained span, so the package carries no omission record.  No retained line contains a citation, so the wrapper carries no bibliography and no bibliography entry.  No retained line contains a figure environment, an image-inclusion command or drawing code, so no image is delivered and the package declares no asset dependency.  Editorial formatting only, in addition: an AMS \texttt{amsart} preamble that restates the article's own declarations and macros used by the retained lines, namely the environments \texttt{theorem} on separate counters, together with the standard packages those lines need.  The retained environments print in this document as Theorem 1 in order of appearance, whereas the article prints the same items with the numbering of its own full document.  The complete native source of the article as distributed in the arXiv e-print for this exact version is bundled unchanged as \texttt{sources/208559/source0.tex}.
\begin{theorem}\label{cov:thm:transversal-matching}
Fix \(t\ge 2\) and \(0<\alpha\le 1/t\). There is \(\gamma=\gamma(t,\alpha)\in(0,1/2)\) such that, for every
\(\kappa>0\), there are \(\xi=\xi(t,\alpha,\kappa)\in(0,1/2)\) and \(N=N(t,\alpha,\kappa)\) with the following property.

Let \(G\) be an \(n\)-vertex graph with \(n\ge N\). Suppose that
\begin{equation*}
V(G)=V_1\dotcupunion\cdots\dotcupunion V_t,
\qquad |V_i|\ge \alpha n\quad\text{for }i\in[t].
\end{equation*}
Write $V_i=U_i\dotcupunion X_i$ and $X=\bigcup_iX_i$. Suppose \(|X|\le \xi n\), and that for \(i\in[t]\),
\begin{equation*}
  \begin{aligned}
    d(u,V_i)&\le \gamma n \qquad&&\text{for }u\in U_i,\\
    d(x,V_j)&\ge \kappa n \qquad&&\text{for }x\in X_i\text{ and }j\ne i.
  \end{aligned}
\end{equation*}
Then
\begin{equation*}
\CC_t(G)\le\prod_{i=1}^t|V_i|.
\end{equation*}
If equality holds, then $G$ is the complete $t$-partite graph with parts $V_1,\ldots,V_t$.
\end{theorem}

\begin{proof}
Suppose we have \(G\) as stated in the theorem, leave \(\gamma, \xi, N\) to be decided. 
Set \(\sigma:=\gamma+\xi<1\). Fix a linear order on every part \(V_i\).
Fix a non-transversal copy of \(K_t\) \(S\), put
\begin{equation*}
r_i:=|S\cap V_i|,
\qquad
A:=\{i:r_i>0\},
\qquad
B:=\{i:r_i=0\},
\qquad
d:=|B|.
\end{equation*}
Let \(P:=(r_1,\cdots,r_t)\) be its multiplicity pattern. Since \(|S|=t\),
\begin{equation}\label{cov:eq:excess-identity}
\sum_{i\in A}(r_i-1)=d.
\end{equation}
For every \(i\in A\), choose a canonical representative \(b_i\in S\cap V_i\):
take the first vertex of \(S\cap U_i\) if this set is nonempty, and otherwise
take the first vertex of \(S\cap X_i\). The other vertices of \(S\cap V_i\) are called \emph{nonrepresentative vertices}.

Given the canonical representatives \(b_i\)'s, the nonrepresentative vertices have at most \((\sigma n)^d\)
possible choices. Indeed, if \(b_i\in U_i\), every nonrepresentative vertex in \(V_i\) lies in \(N(b_i,V_i)\),
which has size at most \(\gamma n\). If \(b_i\in X_i\), then \(S\cap U_i=\varnothing\) by the definition of \(b_i\),
so every nonrepresentative vertex in \(V_i\) lies in \(X_i\), which has size at most \(\xi n\).
And \eqref{cov:eq:excess-identity} says there are exactly \(d\) nonrepresentative vertices to be decided.

We call \(S\) \emph{singular} if \(d=1\) and its unique repeated part \(V_i\) satisfies
\begin{equation*}
S\cap V_i=\{u,x\},
\qquad u\in U_i,
\quad x\in X_i.
\end{equation*}
Let \(V_j\) be the unique missing part. For every \(y\in N(x,V_j)\), form the transversal that
contains \(u\) in \(V_i\), \(y\) in \(V_j\), and the unique vertex of \(S\) in every other part.
This transversal is admissible for \(S\): if it is a clique, its union with \(S\) is a clique;
otherwise it is missing. Thus every singular \(S\) has at least \(\kappa n\) candidates.

A fixed transversal is produced as a candidate by this construction for at most \(t(t-1)|X|\) singular copies.
After the repeated part and the missing part are fixed, the transversal determines \(u\), the vertex chosen
in the missing part, and all singleton vertices of \(S\); only \(x\in X_i\) remains undetermined.

We now define candidates for every nonsingular \(S\). For each completion vector \(y=(y_j)_{j\in B}\) in \(\prod_{j\in B}V_j\), let
\begin{equation*}
T_0(y):=\{b_i:i\in A\}\cup\{y_j:j\in B\}.
\end{equation*}
If \(T_0(y)\) is missing, take \(T_0(y)\) as the candidate. If \(S\cup T_0(y)\) is a clique,
again take \(T_0(y)\). Otherwise both \(S\) and \(T_0(y)\) are cliques but \(S\cup T_0(y)\) is not.
Thus there exist a nonrepresentative vertex \(z\in S\cap V_i\) and a completion vertex \(y_j\) such
that \(zy_j\notin E(G)\). Order the indexed pairs lexicographically using the part indices and the
fixed vertex orders. Choose the first such pair, replace \(b_i\) by \(z\), and take the resulting transversal.
It contains the nonedge \(zy_j\), thus is missing and hence admissible.

Different completion vectors give different candidate transversals.
Thus every nonsingular \(S\) has at least \(\prod_{j\in B}|V_j|\ge \alpha^d n^d\) candidates.

Given a nonsingular multiplicity pattern \(P=(r_1,\ldots,r_t)\) and a transversal \(T\),
if \(T\) serves as \(T_0(y)\) in the above construction,
then \(T\) determines all canonical representatives, leaving at most \((\sigma n)^d\) choices.
Otherwise it replaces some \(b_i\) by a nonrepresentative vertex \(z\), and the pair consisting
of \(i\) and the missing part containing the selected completion vertex has at most \(t(t-1)\) choices.
The transversal \(T\) determines all other representatives.
If \(z\in U_i\), there are at most \(\gamma n\) choices for \(b_i\); if \(z,b_i\in X_i\), there are
at most \(\xi n\); and the remaining nonrepresentative vertices have at most \((\sigma n)^{d-1}\) choices.
In the only remaining case, \(z\in X_i\) and \(b_i\in U_i\). Here there are at most \(n\) choices for \(b_i\),
but nonsingularity implies \(d\ge2\), so this case contributes at most \(\sigma n^d\).
Consequently there is a constant \(C=C(t)\) such that \(T\) is a candidate for at most \(C\sigma n^d\)
nonsingular copies of pattern \(P\) by our construction.

We construct a bipartite \emph{candidate graph} \(H\) as below: its left side is the non-transversal \(K_t\)'s,
grouped by their multiplicity patterns plus one additional pattern for all singular non-transversals;
its right side is \(\mathfrak T\), there is an edge between a non-transversal \(K_t\) \(S\) and a transversal \(T\)
if and only if \(T\) is a candidate for \(S\) obtained through the above construction.
Let \(L=L(t)\) be the number of multiplicity patterns plus one.
Fix \(R\subseteq\mathfrak T\) with \(|R|\le2\). We would like non-transversals of different patterns to use
disjoint classes of candidates. For this purpose, we assign the members of \(\mathfrak T\setminus R\) independently
and uniformly to the \(L\) multiplicity patterns. Later on, for a non-transversal \(K_t\) \(S\) of pattern \(P\),
we only use candidates assigned to \(P\). Specifically, each of \(S\)'s original candidates outside \(R\) is preserved
with probability \(1/L\) independently. Delete the edges between each non-transversal \(K_t\) and
transversals not assigned to its pattern. Denote by \(\mathbf H\) the resulting bipartite graph. Then
\begin{equation*}
\mathbb E|N_{\mathbf H}(S)|=\frac{1}{L}\left|N_{H}(S)\setminus R\right|
\end{equation*}
for any non-transversal \(K_t\) \(S\). Recall that preceding analysis gives
\begin{equation*}
  |N_H(S)|\ge
  \begin{cases}
    \kappa n & S  \mathrm{\ is\ singular},\\
    \alpha^dn^d & S \mathrm{\ is\ nonsingular},
  \end{cases}
\end{equation*}
where \(d\ge1\) when \(S\) is nonsingular. Since \(R\) rules out at most \(2\) candidates, a Chernoff bound gives
\begin{equation*}
  \Pr\left(|N_{\mathbf{H}}(S)|<\frac{|N_H(S)|}{4L}\right)\le e^{-cn}
\end{equation*}
for large enough \(n\), where \(c=\min\{\kappa,\alpha^t\}/(16L)>0\). There are fewer than \(n^t\) non-transversal
\(K_t\)'s and \(n^te^{-cn}<1\) for large enough \(n\). Consequently, there is an assignment such that each
non-transversal \(K_t\) \(S\) retains at least \(\frac{1}{4L}\) of its original candidates.
Denote the resulting bipartite graph under this assignment by \(\widetilde{H}\). By our construction,
\(\widetilde{H}\) is the union of \(L\) disjoint bipartite subgraphs \(\widetilde{H}_P\), one for each pattern,
plus the isolated set \(R\). By preceding analysis, for each transversal \(T\) on the right side assigned to pattern \(P\), we have
\begin{equation*}
  |N_{\widetilde{H}_P}(T)|\le
  \begin{cases}
    t(t-1)\xi n & P  \mathrm{\ is\ the\ singular\ pattern},\\
    C\sigma n^d & P \mathrm{\ is\ nonsingular}.
  \end{cases}
\end{equation*}

Given \(t\ge 2\) and \(0<\alpha\le 1/t\), choose
\begin{equation*}
\gamma:=\min\left\{\frac13,\frac{\alpha^t}{8CL}\right\}.
\end{equation*}
Given \(\kappa>0\), choose \(\xi>0\) so small that
\begin{equation*}
\xi\le\gamma,
\qquad
t(t-1)\xi\le\frac{\kappa}{4L}.
\end{equation*}
Then on every \(\widetilde{H}_P\) we have minimum left degree at least as large as maximum right degree.
Write the two degrees as \(\delta_{\widetilde{H}_P}^{\mathrm{left}}\) and \(\Delta_{\widetilde{H}_P}^{\mathrm{right}}\) respectively.
For every set \(\mathcal S\) of left vertices in \(\widetilde{H}_P\) we have
\begin{equation*}
\delta_{\widetilde{H}_P}^{\mathrm{left}}|\mathcal S|
 \le e(\mathcal S,N(\mathcal S))
 \le \Delta_{\widetilde{H}_P}^{\mathrm{right}}|N(\mathcal S)|,
\end{equation*}
so Hall's condition holds. Combining all the resulting matchings together gives an injection
\begin{equation}\label{cov:eq:transversal-injection}
\phi_R:\mathcal K\longrightarrow\mathfrak T\setminus R,
\end{equation}
where \(\mathcal K\) is the family of all non-transversal copies of \(K_t\) and \(\phi_R(S)\) is admissible for \(S\).

Apply \eqref{cov:eq:transversal-injection} with \(R=\varnothing\).
Start with all actual transversal copies, namely \(\mathcal A\). If \(\phi_\varnothing(S)\) is missing,
add \(S\) to the cover. If \(\phi_\varnothing(S)\in\mathcal A\), replace \(\phi_\varnothing(S)\) by
the clique \(S\cup\phi_\varnothing(S)\). Since \(\phi_\varnothing\) is injective,
at most \(|\mathfrak T\setminus\mathcal A|\) cliques are added for missing transversals and
each actual transversal \(K_t\) is enlarged at most once. Therefore
\begin{equation*}
\CC_t(G)\le |\mathcal A|+|\mathfrak T\setminus\mathcal A|
 =|\mathfrak T|=\prod_{i=1}^t|V_i|.
\end{equation*}

Suppose equality holds. If a missing transversal \(T\) exists, use \(\phi_{\{T\}}\)
from \eqref{cov:eq:transversal-injection}. Then at most \(|\mathfrak T\setminus\mathcal A|-1\)
non-transversal copies are paid for by missing transversals, so the cover has fewer than \(|\mathfrak T|\) members.
Hence every transversal is an actual \(K_t\), and every pair of distinct parts is complete.

If some \(V_i\) contains an edge \(uv\), choose one vertex from every other part. Together with \(u,v\),
these vertices form a \(K_{t+1}\). Let \(T_u,T_v\) be the two transversal copies of \(K_t\) contained in this clique.
Start with \(\mathcal A\setminus\{T_u,T_v\}\), apply \eqref{cov:eq:transversal-injection} with
\(R=\left\{T_u,T_v\right\}\) and add the \(K_{t+1}\) containing \(T_u,T_v\) we just selected.
This gives a cover with at most \(|\mathfrak T|-1\) members, a contradiction.
Thus every part is independent, and \(G\) is complete \(t\)-partite.
\end{proof}

\end{document}
